\documentclass[11pt]{article}
\usepackage[spanish]{babel}
\usepackage[margin=2cm]{geometry}
\usepackage{tabularx}
\providecommand{\toprule}{\hline}
\providecommand{\midrule}{\hline}
\providecommand{\bottomrule}{\hline}
\usepackage[backend=biber,style=ieee]{biblatex}
\addbibresource{referencias.bib}

\begin{document}

\section{Impacto y analisis de la unidad 1}

El incidente de red de Google Cloud del 2 de junio de 2019 afecto a varios servicios y
regiones. Google reporto cifras diferenciadas por servicio, region e intervalo. El impacto no
fue uniforme: algunas regiones sufrieron perdida total de conectividad, mientras que otras
experimentaron degradacion parcial de latencia o errores. \cite{googlecloud19009}

\subsection{Cifras por servicio y region}

En Compute Engine, la region us-east1 presento hasta 33 por ciento de perdida de paquetes
entre las 11:38 y las 12:17. Entre las 12:17 y las 14:50, la perdida se redujo a un maximo de
8 por ciento. En us-central1, la perdida alcanzo un pico de 9 por ciento inmediatamente
despues de las 11:38 y disminuyo hacia las 12:05. \cite{googlecloud19009}

Desde us-west1, la perdida de paquetes hacia us-east1 alcanzo 20 por ciento y hacia
us-central1 alcanzo 8.6 por ciento. Hacia Europa, la perdida fue de hasta 1.9 por ciento. Las
instancias de Compute Engine en us-east4, us-west2, northamerica-northeast1 y
southamerica-east1 quedaron inaccesibles durante el incidente. \cite{googlecloud19009}

Las aplicaciones de App Engine en us-east4, us-west2, northamerica-northeast1 y
southamerica-east1 no estuvieron disponibles durante la interrupcion. En us-central, el
servicio experimento una caida de 23.2 por ciento en solicitudes por segundo. \cite{googlecloud19009}

Cloud Interconnect en us-east4, us-west2, northamerica-northeast1 y southamerica-east1
presento perdida de paquetes entre 50 y 100 por ciento entre las 11:45 y las 16:10. Algunos
enlaces BGP se desconectaron. En us-west1, us-east1 y us-central1, las interconexiones
conectadas a la costa este registraron perdida de paquetes entre 10 y 50 por ciento debido a
la congestion de la red troncal. \cite{googlecloud19009}

Los gateways de Cloud VPN en us-east4, us-west2, northamerica-northeast1 y
southamerica-east1 quedaron inaccesibles. En us-central1 se registro 25 por ciento de perdida
de paquetes y en us-east1 se registro 10 por ciento. \cite{googlecloud19009}

Cloud Pub/Sub presento una caida global de 27 por ciento en publicaciones y suscripciones
exitosas. En us-east4, el error promedio fue de 0.3 por ciento en publicaciones y 25 por
ciento en suscripciones. En southamerica-east1, los errores fueron de 11 por ciento en
publicaciones y 36 por ciento en suscripciones. En northamerica-northeast1, los errores
fueron de 6 por ciento en publicaciones y 31 por ciento en suscripciones. \cite{googlecloud19009}

Cloud Storage presento los mayores porcentajes de error en us-west2, con 96.2 por ciento;
southamerica-east1, con 79.3 por ciento; us-east4, con 62.4 por ciento; y
northamerica-northeast1, con 43.4 por ciento. En us, us-east1, us-west1 y us-central1, el
error promedio fue de 3.5, 1.7, 1.2 y 0.7 por ciento, respectivamente. \cite{googlecloud19009}

BigQuery registro un error global promedio de 0.7 por ciento. La operacion jobs.insert
presento 27 por ciento de errores entre las 11:47 y las 12:02. La operacion
tabledata.insertAll alcanzo un pico de 24 por ciento de errores, aunque su error promedio fue
inferior a 0.01 por ciento. \cite{googlecloud19009}

El frontend de Stackdriver Monitoring experimento una caida de 5 a 10 por ciento en
solicitudes por segundo. El error promedio de inicio de sesion fue de 8.4 por ciento y el
error al cargar datos fue de 3.5 por ciento. \cite{googlecloud19009}

Cloud Endpoints presento errores de hasta 4.4 por ciento al inicio del incidente. Entre las
12:50 y las 15:40, el error promedio fue de 0.6 por ciento. Ademas, el 50 por ciento de los
flujos de publicacion de configuracion fallo entre las 12:00 del 2 de junio y las 11:30 del
4 de junio. \cite{googlecloud19009}

\subsection{Impacto en Cloud Spanner}

Cloud Spanner fue el componente seleccionado para el analisis de CAP. Durante el incidente,
las instancias regionales de Spanner en us-east4, us-west2 y northamerica-northeast1
estuvieron indisponibles entre las 11:48 y las 15:44. \cite{googlecloud19009}

\begin{table}[htbp]
\centering
\caption{Impacto de Cloud Spanner por configuracion y intervalo}
\label{tab:spanner}
\begin{tabular}{p{4cm}p{3.5cm}p{7cm}}
\hline
\textbf{Configuracion} & \textbf{Intervalo} & \textbf{Cifra reportada} \\
\hline
us-east4, us-west2 y northamerica-northeast1 & 11:48 a 15:44 & Instancias regionales
indisponibles. \\
nam3 & 11:50 a 13:10 & Mas de 1 segundo adicional de latencia en el percentil 99. \\
nam3 & 13:10 a 16:50 & 100 ms adicionales de latencia en el percentil 99. \\
nam6 & 11:50 a 13:10 & 320 ms adicionales de latencia en el percentil 99. \\
nam6 & 13:10 a 16:50 & 130 ms adicionales de latencia en el percentil 99. \\
us-central1 & 11:50 a 13:10 & 80 ms adicionales de latencia en el percentil 99. \\
us-east1 & 11:50 a 13:10 & 2 segundos adicionales de latencia en el percentil 99. \\
us-east1 & 13:10 a 15:50 & 250 ms adicionales de latencia en el percentil 99. \\
us-west1 & 11:50 a 14:10 & 20 ms adicionales de latencia. \\
nam-eur-asia1 & 13:50 a 15:20 & 120 ms adicionales de latencia en el percentil 99. \\
\hline
\end{tabular}
\end{table}

Los datos indican que el impacto en Spanner se manifesto de dos formas. En las regiones
us-east4, us-west2 y northamerica-northeast1, el servicio perdio disponibilidad. En otras
configuraciones, el servicio continuo operando, pero con un aumento considerable de latencia.
\cite{googlecloud19009}

\subsection{Transparencia afectada}

El incidente rompio varias formas de transparencia propias de los sistemas distribuidos.

La transparencia de fallo busca ocultar al usuario la ocurrencia de fallos internos. Durante
el incidente, esta transparencia se rompio porque los clientes observaron errores, timeouts,
perdida de paquetes e indisponibilidad de servicios. \cite{googlecloud19009}

La transparencia de ubicacion supone que el usuario no necesita conocer la ubicacion fisica de
los recursos. En este caso, la ubicacion de las regiones afectadas se hizo visible porque los
servicios en us-east4, us-west2, northamerica-northeast1 y southamerica-east1 quedaron
indisponibles. \cite{googlecloud19009}

La transparencia de fallo parcial busca limitar el efecto de un fallo a un componente o region
especifica. Google reporto que la automatizacion desprogramo clusteres de control de red en
varias ubicaciones fisicas al mismo tiempo. Esto amplio el alcance del fallo y evito que la
falla permaneciera contenida en un solo dominio. \cite{googlecloud19009}

La transparencia de replicacion busca que el usuario no perciba la existencia de replicas o
copias distribuidas de los datos. En Cloud Spanner, la replicacion no pudo ocultar
completamente la interrupcion. Las instancias regionales afectadas dejaron de estar
disponibles y las configuraciones multiregion experimentaron aumentos de latencia.
\cite{googlecloud19009}

\subsection{Aplicacion de CAP a Cloud Spanner}

El teorema CAP establece que, ante una particion de red, un sistema distribuido no puede
garantizar simultaneamente consistencia y disponibilidad. En este caso, la congestion de la
red troncal y la perdida de rutas BGP constituyen una degradacion severa de la red que puede
analizarse como una particion. \cite{googlecloud19009}

\begin{table}[htbp]
\centering
\caption{Aplicacion de CAP a Cloud Spanner}
\label{tab:cap}
\begin{tabular}{p{3cm}p{11.5cm}}
\hline
\textbf{Propiedad} & \textbf{Aplicacion al incidente} \\
\hline
Consistencia & Spanner mantiene transacciones distribuidas externamente consistentes. El
sistema no expone datos que violen el orden de las transacciones. \cite{corbett2013spanner} \\
Disponibilidad & Las instancias regionales de Spanner en us-east4, us-west2 y
northamerica-northeast1 estuvieron indisponibles entre las 11:48 y las 15:44. El servicio
prefirio no responder antes que comprometer la consistencia. \cite{googlecloud19009} \\
Tolerancia a particiones & La perdida de rutas BGP y la congestion del backbone actuaron como
una particion o degradacion severa de la red. Spanner continuo operando bajo esa condicion en
las configuraciones que conservaron conectividad. \cite{googlecloud19009} \\
\hline
\end{tabular}
\end{table}

La evidencia del incidente muestra que Spanner priorizo la consistencia. Ante la degradacion
de la red, el sistema redujo disponibilidad en las regiones afectadas y aumento la latencia en
otras configuraciones, en lugar de ofrecer respuestas que pudieran comprometer la consistencia
externa. \cite{googlecloud19009,corbett2013spanner}

\subsection{Mecanismo de tolerancia a fallos insuficiente}

El mecanismo de tolerancia a fallos que resulto insuficiente fue la redundancia del plano de
control de red. Google disponia de multiples instancias del software de gestion de clusteres,
pero una configuracion las marco como elegibles para el mismo evento de mantenimiento.
\cite{googlecloud19009}

El software de mantenimiento tenia un error que permitia desprogramar varios clusteres
independientes al mismo tiempo, incluso cuando estaban ubicados en sitios fisicos distintos.
Esto violo el principio de aislamiento entre dominios de fallo. \cite{googlecloud19009}

La red estaba disenada para operar temporalmente sin plano de control en un modo conocido como
fail static. Sin embargo, ese periodo fue demasiado corto. Despues de varios minutos, se
retiraron rutas BGP entre ubicaciones afectadas, lo que redujo la capacidad de la red troncal
y desconecto regiones. \cite{googlecloud19009}

La recuperacion tambien fue lenta porque la configuracion del plano de control se perdio y
tuvo que reconstruirse y redistribuirse. Google indico que persistir dicha configuracion
habria reducido el tiempo de recuperacion en un orden de magnitud. \cite{googlecloud19009}

\printbibliography

\end{document}
